\SourceLeaf{393}%
straight line of course fails. In this case, however, it is easy to show that the common midpoint must also be the centre of the ellipse.

2) If we extend two opposite sides of the quadrilateral until they meet, and do the same for the other two, we may also regard the straight segment between the two points of intersection as a vanishing ellipse tangent to all four sides of the quadrilateral. Its midpoint must therefore lie on the same straight line that joins the midpoints of the two diagonals. To my knowledge, this general property of every quadrilateral has yet to be noted. I shall give a simple direct proof below.

To shorten the calculations still further, I shall now suppose that this straight line itself has been chosen as the axis of abscissas, so that \(\varphi=0\). The origin of the abscissas remains arbitrary as before. This very choice of \(\varphi=0\) makes one of the four fundamental equations unnecessary. To determine the four unknown quantities \(t,u,r,\psi\), we thus have the three equations
\[
\begin{aligned}
2aa+t-4ar\cos A+rr\cos2A-u\cos2(A-\psi)&=0\\
2a'a'+t-4a'r\cos A'+rr\cos2A'-u\cos2(A'-\psi)&=0\\
2a''a''+t-4a''r\cos A''+rr\cos2A''-u\cos2(A''-\psi)&=0
\end{aligned}
\]
together with the condition that the area of the ellipse, to which the product \(\alpha\beta\) is clearly proportional. and therefore also \(4\alpha\alpha\beta\beta\) or \((rr-t)^2-uu\), is to be a maximum.

If we set \(rr-t=\theta\) for brevity, and
\[
\begin{aligned}
b&=2(a-r\cos A)^2\\
b'&=2(a'-r\cos A')^2\\
b''&=2(a''-r\cos A'')^2
\end{aligned}
\]
then the equations above become
\[
\begin{aligned}
\theta+u\cos2(A-\psi)&=b\\
\theta+u\cos2(A'-\psi)&=b'\\
\theta+u\cos2(A''-\psi)&=b''
\end{aligned}
\]
